\section{Related Work}

\paragraph{Agents in Finance} 
Recent advances in financial LLMs \citep{liu2023fingpt, liu2025fin} and evaluation benchmarks \citep{guo2025fineval, tang2025finmmr} have substantially expanded the scope of automated financial reasoning.
Building on these foundations, a growing line of work explores how agentic systems \citep{li2025investorbench, yang2026finvault} can move beyond financial understanding toward {decision-oriented research workflows}, including factor discovery and trading analysis.
In this context, the pursuit of alpha factors has evolved from manual engineering and heuristic search toward LLM-driven closed-loop discovery \cite{li2024can,shi2025navigating, chen2025alphasage, wang2025alpha}.
Beyond search optimization, recent efforts emphasize workflow systematization. RD-Agent \cite{li2025rdagent} proposes a multi-agent framework that decouples the pipeline into research and development stages, enabling data-centric joint optimization of factors and models. To address market non-stationarity, AlphaForge \cite{shi2024alphaforge} focuses on the dynamic combination of mined factors, 
while Alphafin \cite{li2024alphafin} and AlphaEval \cite{ding2025alphaeval} establish standardized, task-oriented protocols for reproducible evaluation. Despite these advances, trading constraints (e.g., turnover, complexity) remain largely post-hoc filters rather than intrinsic objectives, leading to suboptimal generalization and interpretability in live trading.
\paragraph{Self-Evolving Agents} Self-evolving agents \cite{fang2025comprehensive, zhai2025agentevolver, lin2025se, zhang2026evofsm} represent a paradigm shift from static instruction-following to autonomous learning through interacting with environment. 
AlphaEvolve \cite{novikov2025alphaevolve} demonstrates a coding-centric evolution where the agent employs evolutionary resampling to autonomously generate algorithms for scientific discovery. Building on this, CSE \cite{hu2026controlled} introduces controllable self-evolution, facilitating a critical transition from stochastic generation to feedback-driven evolution.
The financial domain adapts thi s self-evolutionary framework to handle non-stationary and high-dimensional data. Research has converged on multi-agent coordination and structured feedback to guide evolution \cite{li2025hedgeagents}. TradingAgents \cite{xiao2025trad} and FactorMAD \cite{duan2025factormad} utilize institutional-style debates to refine trading hypotheses, while QuantAgents \cite{li2025quantagents} and ATLAS \cite{papadakis2025atlas} incorporate simulated trading performance as a reward signal for dynamic prompt optimization. To ensure consistency over long horizons, FinMem \cite{yu2025finmem} and FinCon \cite{yu2024fincon} leverage hierarchical memory and conceptual reinforcement to retain and refine high-level trading experience.
Despite these advancements, the transition of self-evolving agents to finance is hindered by the low signal-to-noise ratio and non-stationary. To mitigate this, our \method employs stable optimization units with mutation and crossover operators, ensuring a robust, traceable evolution path through structured archiving of iteration logic.
  